\section{GitHub Actions Workflow}

With the containers working locally, the next step was to automate the build and push
process. I set up a GitHub Actions workflow so that I never have to run
\texttt{podman build} manually again. The workflow is defined in
\texttt{.github/workflows/ci.yml} and covers testing, building multi-architecture
images, scanning for vulnerabilities, and sending a notification when it is done.

\subsection{Triggers}

The workflow runs on two events: a \texttt{workflow\_dispatch}, which lets me trigger it
manually from the GitHub UI, and any pull request targeting the main branch. I chose not
to trigger it on every push to main. The idea is that the pipeline validates a pull
request before it is merged, and I can also kick it off by hand whenever needed.

\subsection{Job 1: Backend Tests}

The first job checks out the code, installs Go 1.24 with dependency caching pointed at
\texttt{backend/go.sum}, and runs \texttt{make test} inside the \texttt{backend/}
directory. If the tests fail the rest of the pipeline stops. There is no point building
and pushing a container image if the code inside it is already broken.

\subsection{Job 2: Build and Push}

The build job uses a matrix strategy so it runs once for the backend and once for the
frontend in parallel, rather than having to write two separate near-identical jobs:

\begin{lstlisting}[language=yaml]
strategy:
  matrix:
    component: [backend, frontend]
\end{lstlisting}

Inside the job I set up QEMU and Docker Buildx to enable cross-platform builds. The
images are built for both \texttt{linux/amd64} and \texttt{linux/arm64} so they work on
standard cloud servers as well as on my ARM64 Mac when testing locally.

For tagging I used \texttt{docker/metadata-action} rather than constructing tags by
hand. It generates three tags automatically: one based on the commit SHA
(\texttt{sha-xxxxxxx}), one based on the branch name, and \texttt{latest} when the
build is on the default branch:

\begin{lstlisting}[language=yaml]
tags: |
  type=sha
  type=ref,event=branch
  type=raw,value=latest,enable={{is_default_branch}}
\end{lstlisting}

The SHA tag is the most useful in practice because it ties every image back to an exact
commit, which makes debugging a bad deployment much easier. Images are only pushed when
the event is not a pull request, so branches that have not been reviewed do not pollute
GHCR with unfinished work.

After the push, the job runs Trivy to scan the published image. I configured it to exit
with a failure code if it finds \texttt{CRITICAL} or \texttt{HIGH} severity issues that
have a fix available. Unfixed vulnerabilities are ignored since there is nothing I can
act on immediately:

\begin{lstlisting}[language=yaml]
- name: Run Trivy vulnerability scanner
  uses: aquasecurity/trivy-action@master
  with:
    image-ref: >-
      ${{ env.REGISTRY }}/${{ env.IMAGE_NAME }}-${{ matrix.component }}:latest
    format: "table"
    exit-code: "1"
    ignore-unfixed: true
    vuln-type: "os,library"
    severity: "CRITICAL,HIGH"
\end{lstlisting}

\subsection{Job 3: Discord Notification}

The last job runs after both the test and build jobs have finished, regardless of whether
they succeeded or failed, using \texttt{if: always()}. It sends a message to a Discord
channel with the overall build status, the repository name, the event that triggered the
run, and a direct link to the workflow on GitHub. The webhook credentials are stored as
repository secrets.

This turned out to be more useful than I expected. Instead of having to check GitHub
constantly, I get a notification in Discord the moment a build completes, and the status
line tells me immediately whether it passed or failed.

\subsection{One Problem with Actions}

One error I ran into early was a confusing one. The workflow failed with:

\begin{lstlisting}
Error: Unable to resolve action actions/setup-buildx-action,
repository not found.
\end{lstlisting}

I had assumed all the standard actions were under the \texttt{actions/} namespace, the
same as \texttt{actions/checkout} and \texttt{actions/setup-go}. It turned out that the
Buildx and QEMU actions are maintained by Docker and live under the \texttt{docker/}
namespace. The fix was to change the paths to \texttt{docker/setup-qemu-action} and
\texttt{docker/setup-buildx-action}. It is the kind of mistake you only make once.

\subsection{What the Pipeline Does Not Do Yet}

The current workflow does not automatically update the infrastructure repository with the
new image tag, which is the step that would close the GitOps loop and let ArgoCD pick
up new versions on its own. That requires a separate job that checks out the config
repo, uses \texttt{yq} to patch the image reference in the deployment manifest, and
pushes the commit using a personal access token. I built and tested this in a separate
branch during the ArgoCD work described in section 5, but kept it out of the main
workflow since the ArgoCD integration itself could not be completed on the Sandbox. The
wiring is there and ready; it just needs the right cluster environment to be worth
enabling.

\subsection{Two Repositories and Why}

The application code and the OpenShift manifests live in separate Git repositories. The
application repo contains the source code, Containerfiles, and this CI workflow. The
config repo contains only the Kubernetes YAML files under a \texttt{k8s/} directory.
This follows the GitOps principle that the desired state of the infrastructure should be
its own versioned artifact, separate from the application that runs in it.

The reason I find most compelling is access control. I can give someone access to the
application repo without giving them any ability to change what is deployed to the
cluster. Changes to the cluster go through the config repo, which can have its own
review process. There is also a clean audit trail: if something breaks in production,
the config repo's commit history shows exactly what changed and when, independently of
any application code changes happening at the same time.
